\documentclass{article}
\usepackage[T1]{fontenc}
\usepackage{lmodern}
\usepackage{graphicx}
\usepackage{amsmath,amssymb}
\usepackage[a4paper,margin=2cm]{geometry}
\usepackage[absolute,overlay]{textpos}
\setlength{\TPHorizModule}{1mm}
\setlength{\TPVertModule}{1mm}
\usepackage[indonesian]{babel}
\usepackage{hyperref}
\setlength{\emergencystretch}{2em}
% unit: Halaman 49

\begin{document}

% === Halaman 49 (llm) ===

\section*{Program enkripsi-dekripsi AES dengan Python}
Input: File biner

\vspace{10mm}

\begin{verbatim}
# Enkripsi file biner dengan AES-128
from Crypto.Cipher import AES
from Crypto.Random import get_random_bytes
from Crypto.Util.Padding import pad, unpad
import base64

BLOCK_SIZE = 16           # Panjang blok pesan yang dienkripsi = 16 byte = 128 bit

def enkripsi_aes(plainteks, kunci):
    kunci = kunci.encode()                # Kodekan kunci dari string menjadi byte
    iv = get_random_bytes(BLOCK_SIZE)     # Bangkitkan IV
    AEScipher = AES.new(kunci, AES.MODE_CBC, iv) # Definisikan AES mode CBC

    plainteks = pad(plainteks, BLOCK_SIZE) # lakukan padding jika ukuran plainteks bukan kelipatan 16 byte

    cipherteks = AEScipher.encrypt(plainteks) # Lakukan enkripsi plainteks
    cipherteks = base64.b64encode(iv+cipherteks) # Kodekan iv+cipherteks dari kode byte ke base64
    return cipherteks
\end{verbatim}

\end{document}
